\chapter{CONCLUSÃO}
\label{cap:conclusao}

Este trabalho investigou técnicas computacionais aplicadas ao ensino de matemática por meio de uma revisão sistemática da literatura e da elaboração de uma especificação conceitual de protótipo. O processo organizou registros provenientes de quatro fontes científicas e reteve operacionalmente 18 registros: 17 candidatos empíricos provisórios e 1 protocolo ou proposta contextual. A síntese empírica preliminar apontou recorrência de abordagens voltadas à predição de desempenho e à estimativa de proficiência.

A análise evidenciou que resultados obtidos em bases específicas não podem ser generalizados automaticamente e que indicadores de desempenho não equivalem, isoladamente, à aprendizagem. A especificação também distinguiu a probabilidade preditiva de um alvo definido, entendida como saída condicionada aos dados e ao modelo, do indicador de desempenho observado; essa probabilidade não equivale, isoladamente, à proficiência, à competência, à aprendizagem ou à eficácia pedagógica. A apreciação metodológica preliminar pelo MMAT foi organizada por critérios, sem pontuação geral, preservando a visibilidade das limitações de cada estudo e a rastreabilidade dos julgamentos realizados.

A fundamentação pedagógica permitiu distinguir desempenho observado, proficiência estimada, competência e aprendizagem. Essa diferenciação orientou a especificação do protótipo e reforçou que modelos computacionais devem apoiar a interpretação docente, sem produzir diagnósticos definitivos ou substituir decisões pedagógicas contextualizadas.

No escopo executado, o objetivo geral foi alcançado ao mapear as principais aplicações computacionais e converter as evidências encontradas em uma especificação técnica e pedagógica. Foram definidos princípios de projeto, requisitos funcionais e não funcionais, critérios para fontes de dados e modelos, um protocolo de avaliação e uma arquitetura de referência.

\section{Contribuições do Trabalho}

Entre as contribuições, destaca-se a síntese estruturada da literatura, apoiada por artefatos versionados para reprodução do snapshot e por um \textit{pipeline} automatizado para coleta, deduplicação, triagem e exportação dos registros. O trabalho também separou explicitamente relevância temática e avaliação metodológica, preservando a apreciação preliminar dos estudos empíricos pelo MMAT em respostas organizadas por critério.

A fundamentação sobre aprendizagem matemática e os limites das inferências computacionais orientou uma especificação técnica e pedagógica de protótipo baseada nas evidências da revisão. Essa especificação definiu critérios de rastreabilidade, explicabilidade, avaliação e participação docente, sem transformar uma proposta conceitual em evidência de eficácia pedagógica.

\section{Limitações}

A revisão foi conduzida por um único pesquisador, sem dupla seleção e dupla avaliação independentes. Parte dos julgamentos dependeu de resumos, metadados e descrições metodológicas disponíveis, o que restringiu a profundidade da apreciação de alguns estudos. A recuperação de fontes, os localizadores dos critérios e a adjudicação do MMAT ainda precisam ser concluídos antes de qualquer síntese final de qualidade. As fontes consultadas e os idiomas selecionados também limitaram a cobertura da literatura.

A especificação do protótipo não incluiu uma aplicação funcional, o treinamento de modelos com uma base definitiva ou validação com participantes. Consequentemente, o trabalho não produziu métricas próprias de acurácia nem evidências de eficácia pedagógica. Essa limitação foi assumida como delimitação do escopo executado; a contribuição técnica ficou delimitada à engenharia de requisitos e à arquitetura conceitual.

\section{Continuidade da Pesquisa}

Como continuidade, recomenda-se implementar a especificação em um ambiente reproduzível, selecionar uma base de dados compatível com os critérios definidos e comparar modelos de referência sob o protocolo de avaliação proposto. Estudos com participantes devem observar autorização institucional, proteção dos dados e participação de profissionais da educação na interpretação dos resultados.

A continuidade também deve investigar estabilidade das explicações, diferenças de desempenho entre grupos, alinhamento com descritores curriculares e utilidade das informações para o planejamento docente. Essas atividades ampliam os resultados obtidos, mas não alteram a conclusão de que a revisão e a especificação responderam ao escopo definido para este trabalho.

%--------- FIM CONCLUSÃO ------------
